\section{总结与展望}\label{sec:conclusion}

\subsection{工作总结}

本文围绕遥感图像中的微小船舶检测，完成了文献梳理、模型实验和系统实现几项工作。研究部分先整理通用目标检测、小目标检测和遥感船舶检测的发展脉络，再在同一数据和评测口径下比较 DRENet、FCOS 与 YOLO26。工程部分采用 React、FastAPI 和 SQLite 搭建本地检测系统，完成任务提交、异步推理、结果可视化和历史回看。

实验部分的主要结论是，DRENet 和 YOLO26 的 AP50 非常接近，但二者取向不同。DRENet 的 Recall 更高，适合先尽量发现目标的场景；YOLO26 在 AP50-95、Precision 和 F1 上更均衡，更适合作为系统默认模型。FCOS 的整体结果弱于前两者，但保留了 anchor-free 路线的参照意义。阈值消融和输入尺寸敏感性分析也说明，误检、漏检和尺度扰动不能只靠单个 AP50 指标解释。

系统部分并不是单独做一个结果展示页，而是把推理能力、任务状态、结果文件和历史记录放到同一套流程里管理。经过联调后，系统已经可以接入真实模型权重，支持任务回放和答辩演示，离线实验中的部分结果也能在页面中复查。

本文的工作重点不在提出新的检测网络，而在于把数据口径、实验记录和系统实现整理成一条可复查的路径。对本科毕业设计而言，这条路径可以支撑论文写作，也可以支撑现场演示。

\subsection{不足分析}

当前工作已经完成主要目标，但仍有几处边界需要说明。

第一，输入尺寸敏感性分析只覆盖 YOLO26 与 FCOS。DRENet 在本地插件路径下使用非 512 输入会触发 shape mismatch，因此没有进入统一尺寸主表。第二，FCOS 的正式结果来自 MMDetection 默认输入策略，训练口径和效率统计方式与 DRENet、YOLO26 不完全一致，更适合作为参考模型。第三，系统虽然已能运行，但插件化部署规范和跨机器迁移说明还可以继续补充，部分模型接入仍依赖本地代码路径。第四，融合模式已经跑通，但还没有形成统一离线 AP 主表，所以未作为核心结论展开。

这些问题更多属于后续完善空间，不改变本文对 DRENet、YOLO26 和 FCOS 的基本判断。若继续提高可迁移性和评测严谨性，重点应放在部署规范、统一离线评估和跨模型对齐上。

\subsection{未来展望}

后续可以从以下方向继续推进。

一是补齐融合模式的统一离线评测，把单模型结果和融合展示之间的关系量化出来。二是继续围绕 DRENet 与 YOLO26 的误检、漏检取舍调整阈值、后处理和融合策略。三是重新统一 FCOS 的输入尺度与训练策略，使跨模型比较更严谨。四是整理模型插件化接入和部署说明，降低换机器运行时的成本。五是在更多遥感场景和数据集上检查方法与系统的泛化情况。

沿这些方向继续做下去，现有系统可以从演示平台逐步扩展为实验管理平台。例如，把固定演示样例、统一报告导出、批量回评和参数留痕接入同一工作流，让实验复查和系统展示共享更多底层能力。

\subsection{全文结论}

本文完成了遥感微小船舶检测任务中的实验比较和系统实现。在 LEVIR-Ship 数据集上，DRENet 更偏召回优先，YOLO26 更偏精度和部署均衡，FCOS 保留为 anchor-free 参照。系统实现则说明，将推理能力服务化并纳入任务管理后，结果展示、历史回放和演示交付会更加清楚。

从最终完成情况看，本文已经把问题分析、实验论证和系统实现串联起来。现有实验记录、分析框架和系统代码，为后续性能优化、评测补齐和部署规范化留下了可继续迭代的基础。
